\documentclass{beamer}
\usepackage{risk231}   % shared preamble: packages, theorems, colours, listings

\setchapter{5}         % sections numbered 5.1, 5.2, ...

\title{RISK 231: SQL}

%----------------------------------------------------------------------
% Local helpers for this chapter
%----------------------------------------------------------------------
\definecolor{joinfill}{RGB}{237,160,145}
% circled step number for the SQL order of execution
\newcommand{\stepnum}[1]{\tikz[baseline=(c.base)]\node[circle,draw,thick,inner sep=1pt,
  minimum size=1.2em,font=\scriptsize\bfseries](c){#1};}
% Venn diagram for a join: \joinvenn{title}{A-only fill}{overlap fill}{B-only fill}
\newcommand{\joinvenn}[4]{%
\begin{tikzpicture}[scale=0.7]
  \fill[#2] (0,0) circle (1);
  \fill[#4] (1.3,0) circle (1);
  \begin{scope}\clip (0,0) circle (1); \fill[#3] (1.3,0) circle (1);\end{scope}
  \draw (0,0) circle (1) (1.3,0) circle (1);
  \node[font=\tiny] at (-0.4,0) {Table A};
  \node[font=\tiny] at (1.7,0) {Table B};
  \node[font=\footnotesize\bfseries] at (0.65,1.3) {#1};
\end{tikzpicture}}

\begin{document}
\maketitle

%======================================================================
\section{Recap of Codecademy and Joins}
%======================================================================

\sectionframe

\begin{frame}{Review: Manipulation}
\begin{itemize}
  \item \textbf{SQL} is a programming language designed to manipulate and manage data
        stored in relational databases.
  \begin{itemize}
    \item A \textbf{relational database} is a database that organizes information into
          one or more tables.
    \item A \textbf{table} is a collection of data organized into rows and columns.
  \end{itemize}
\end{itemize}
\vspace{0.5em}
A \textbf{statement} is a string of characters that the database recognizes as a valid command.
\begin{itemize}
  \item \code{CREATE TABLE} creates a new table.
  \item \code{INSERT INTO} adds a new row to a table.
  \item \code{SELECT} queries data from a table.
  \item \code{ALTER TABLE} changes an existing table.
  \item \code{UPDATE} edits a row in a table.
  \item \code{DELETE FROM} deletes rows from a table.
\end{itemize}
\end{frame}

\begin{frame}{Review: Queries}
\begin{itemize}
  \item \code{SELECT} is the clause we use every time we want to query information from a database.
  \item \code{AS} renames a column or table.
  \item \code{DISTINCT} returns unique values.
  \item \code{WHERE} is a popular command that lets you filter the results of the query
        based on conditions that you specify.
  \item \code{LIKE} and \code{BETWEEN} are special operators that work with numbers,
        text, or date data types.
  \item \code{AND} and \code{OR} combine multiple conditions.
  \item \code{ORDER BY} sorts the result.
  \item \code{LIMIT} specifies the maximum number of rows that the query will return.
  \item \code{CASE} creates different outputs.
\end{itemize}
\end{frame}

\begin{frame}{Review: Aggregate functions}
\begin{itemize}
  \item \code{COUNT()}: count the number of rows
  \item \code{SUM()}: the sum of the values in a column
  \item \code{MAX()} / \code{MIN()}: the largest/smallest value
  \item \code{AVG()}: the average of the values in a column
  \item \code{ROUND()}: round the values in the column
\end{itemize}
\vspace{0.5em}
\textbf{Aggregate functions} combine multiple rows together to form a single value of
more meaningful information.
\begin{itemize}
  \item \code{GROUP BY} is a clause used with aggregate functions to combine data from
        one or more columns.
  \item \code{HAVING} limits the results of a query based on an aggregate property.
\end{itemize}
\end{frame}

\begin{frame}{Review: Multiple tables}
\begin{itemize}\small
  \item \code{JOIN} will combine rows from different tables if the join condition is true.
  \item \code{LEFT JOIN} will return every row in the left table, and if the join condition
        is not met, \code{NULL} values are used to fill in the columns from the right table.
  \item \textbf{Primary key} is a column that serves as a unique identifier for the rows in the table.
  \item \textbf{Foreign key} is a column that contains the primary key of another table.
  \item \code{CROSS JOIN} lets us combine all rows of one table with all rows of another table.
  \item \code{UNION} stacks one dataset on top of another.
  \item \code{WITH} allows us to define one or more temporary tables that can be used in the final query.
\end{itemize}
\end{frame}

\begin{frame}{Databases \& Schemas}
\begin{columns}[T]
\begin{column}{0.6\textwidth}
\begin{itemize}\small
  \item \textbf{Database:} an organized collection of data.
  \item \textbf{Schema:} the organization of rule sets that define the structure and
        relationship of data within a database.
  \item Within a database you may have multiple schemas; schemas can have multiple
        \textbf{tables} (datasets of a particular type or for a particular purpose).
  \item \textbf{Views} (\code{\_VW}) are easy ways to execute a commonly used query.
\end{itemize}
\end{column}
\begin{column}{0.35\textwidth}
\centering
\includegraphics[width=0.85\textwidth]{ch05/filing_cabinets}
\end{column}
\end{columns}
\vspace{0.3em}
{\small\textit{For example:}}\par
\centering
\begin{tikzpicture}[font=\scriptsize\itshape, every node/.style={draw=blue!50!black,
    text=blue!50!black, inner sep=2pt}, ->, >=Stealth, red!70!black]
  \node[font=\footnotesize\bfseries\itshape, text width=1.5cm] (db) at (0,0) {Beatles Database};
  \node[anchor=west] (j) at (2.2, 0.75) {John's Schema};
  \node[anchor=west] (p) at (2.2, 0.25) {Paul's Schema};
  \node[anchor=west] (r) at (2.2,-0.25) {Ringo's Schema};
  \node[anchor=west] (g) at (2.2,-0.75) {George's Schema};
  \foreach \s in {j,p,r,g} \draw (db.east) -- (\s.west);
  \node[anchor=west, align=left, text width=3.6cm] (t) at (5.3,0.15)
    {$\bullet$ Ticket\_to\_Ride\_Table\\ $\bullet$ Help\_Table\\ $\bullet$ Norwegian\_Wood\_Table\\
     $\bullet$ I\_Am\_the\_Walrus\_Table\\[2pt] $\bullet$ Come\_Together\_VW};
  \foreach \y in {0.85,0.55,0.25,-0.05,-0.4} \draw (j.east) -- (5.3,\y);
\end{tikzpicture}
\end{frame}

\begin{frame}{SQL Editor}
\centering
\vfill
{\large\color{blue!50!black}\textbf{Let's take a look at the SQL App}}
\vfill
\end{frame}

\begin{frame}{SQL Order of Execution}
Unlike some query languages, SQL does not execute its statements top-down; the clauses
are executed in a particular order:
\vspace{0.6em}
\begin{center}
\renewcommand{\arraystretch}{1.25}
\begin{tabular}{@{}l l@{}}
  \code{SELECT} \stepnum{5} \ \code{DISTINCT} \dots & \stepnum{6} \\
  \code{FROM} \dots & \stepnum{1} \\
  \code{JOIN} \dots\ \code{ON} \dots & \stepnum{1} \\
  \code{WHERE} \dots & \stepnum{2} \\
  \code{GROUP BY} \dots & \stepnum{3} \\
  \code{HAVING} \dots & \stepnum{4} \\
  \code{ORDER BY} \dots & \stepnum{7} \\
  \code{LIMIT}/\code{TOP} {\scriptsize(PL/SQL: \code{WHERE ROWNUM <= \#})} \dots & \stepnum{8} \\
  \code{;} &
\end{tabular}
\end{center}
\end{frame}

\begin{frame}[fragile]{SQL Commenting}
Two ways to comment in SQL:
\begin{itemize}
  \item Double dashes at the front of the line: \lstinline[language=SQL]|-- comment|
  \item Forward slash asterisk on each side of the comment: \lstinline[language=SQL]|/* comment */|
\end{itemize}
\vspace{0.5em}
For example:
\begin{lstlisting}[language=SQL]
--Code to help look up some tables in the database

SELECT * FROM IRIS   /*Special table*/
 WHERE ID > 477156   /*Filters the data*/

--Cool, right?
\end{lstlisting}
\end{frame}

\begin{frame}{SQL Recap: Practice}
Answer the following by writing SQL code in the SQL App:
\begin{prob}
\begin{enumerate}
  \item Using the \code{COUNT()} function, determine how many rows are in the
        \code{SCHEMA\_A.iris} table.
  \item From the \code{SCHEMA\_A.mtcars} table, select all rows that are \underline{not}
        \code{cyl} 6 and have an \code{hp} greater than or equal to 105.
  \item Using the \code{SUM()} function, display the sum of \code{len} for each level of
        \code{supp} in the \code{SCHEMA\_B.ToothGrowth} table.
\end{enumerate}
\end{prob}
\end{frame}

\begin{frame}[fragile]{SQL Recap: Practice -- \red{Answers}}
\small
1. Number of rows in the iris table: \red{\textbf{150}}
\begin{lstlisting}[language=SQL]
SELECT COUNT(*) FROM SCHEMA_A.IRIS
\end{lstlisting}
2. Rows of mtcars that are not \code{cyl} 6 and have \code{hp} $\ge$ 105:
\begin{lstlisting}[language=SQL]
SELECT * FROM SCHEMA_A.MTCARS
WHERE CYL <> 6
AND HP >= 105
\end{lstlisting}
3. Sum of \code{len} for each level of \code{supp}: \red{\textbf{OJ = 619.9 and VC = 508.9}}
\begin{lstlisting}[language=SQL]
SELECT SUPP, SUM(LEN) FROM SCHEMA_B.TOOTHGROWTH
GROUP BY SUPP
\end{lstlisting}
\end{frame}

\begin{frame}[fragile]{Table Management: Creating and Dropping Tables}
\begin{itemize}
  \item \code{CREATE TABLE} -- constructs a table in a given schema and database
\end{itemize}
\begin{lstlisting}[language=SQL]
CREATE TABLE customers (customer_id number(10),
                        customer_name varchar2(50),
CONSTRAINT customers_pk PRIMARY KEY (customer_id));
\end{lstlisting}
\begin{itemize}
  \item \code{CREATE TABLE AS} -- constructs a table from a query
\end{itemize}
\begin{lstlisting}[language=SQL]
CREATE TABLE customers AS
    (SELECT * FROM SCHEMA.TABLE)
\end{lstlisting}
\begin{itemize}
  \item \code{DROP TABLE} -- removes a table in a given schema and database
\end{itemize}
\begin{lstlisting}[language=SQL]
DROP TABLE customers;
\end{lstlisting}
\end{frame}

\begin{frame}[fragile]{Table Management: Adding Columns and Rows}
\begin{itemize}
  \item \code{ALTER TABLE} -- add, modify, or drop/delete columns in a table
\end{itemize}
\begin{lstlisting}[language=SQL]
ALTER TABLE customers ADD (city varchar2(40)
                           DEFAULT 'Seattle');
\end{lstlisting}
\vspace{0.5em}
\begin{itemize}
  \item \code{INSERT INTO} -- insert a single record or multiple records into a table
\end{itemize}
\begin{lstlisting}[language=SQL]
INSERT INTO customers
(customer_id, customer_name)
VALUES
(123456, 'HUSKY');
\end{lstlisting}
\end{frame}

\begin{frame}{Table Management: Joining Tables}
\centering\setlength{\tabcolsep}{2pt}
\begin{tabular}{cccc}
\joinvenn{Inner Join}{white}{joinfill}{white} &
\joinvenn{Left Join}{joinfill}{joinfill}{white} &
\joinvenn{Right Join}{white}{joinfill}{joinfill} &
\joinvenn{Full Join}{joinfill}{joinfill}{joinfill} \\[2pt]
\parbox[t]{0.22\textwidth}{\raggedright\scriptsize Select all records from Table A and Table B,
  where the join condition is met.} &
\parbox[t]{0.22\textwidth}{\raggedright\scriptsize Select all records from Table A, along with records
  from Table B for which the join condition is met (if at all).} &
\parbox[t]{0.22\textwidth}{\raggedright\scriptsize Select all records from Table B, along with records
  from Table A for which the join condition is met (if at all).} &
\parbox[t]{0.22\textwidth}{\raggedright\scriptsize Select all records from Table A and Table B,
  regardless of whether the join condition is met or not.}
\end{tabular}
\end{frame}

\begin{frame}{Joining Tables: Example Data}
Throughout the join examples we use the following two tables:
\vspace{0.6em}
\begin{center}
{\scriptsize\setlength{\tabcolsep}{4pt}
\begin{tabular}[t]{@{}lccc@{}}
\multicolumn{4}{@{}l}{\textbf{Students}}\\ \toprule
Student\_Name & Grad\_Year & Class & Grade \\ \midrule
Patrick & 9 & Lunch & A+ \\
Sandy & 11 & Gym & A \\
Gary & 9 & Civics & B+ \\
Pearl & 12 & Calculus & B \\
Karen & 12 & Calculus & C- \\
Larry & 10 & Spelling & B \\
Plankton & 11 & Gym & B- \\
Squidward & 10 & Biology & A- \\
Bob & 11 & Biology & B+ \\ \bottomrule
\end{tabular}\hspace{2.5em}
\begin{tabular}[t]{@{}ll@{}}
\multicolumn{2}{@{}l}{\textbf{Teachers}}\\ \toprule
Teacher\_Name & Class \\ \midrule
N. Watterson & Calculus \\
B. Joe & Gym \\
T. Wilson & Spelling \\
C. Krueger & Biology \\
S. Small & Civics \\
R. Robinson & Drama \\
B. Parham & Band \\
G. Watterson & Music Theory \\ \bottomrule
\end{tabular}}
\end{center}
\end{frame}

\begin{frame}[fragile]{Joining Tables: \code{INNER JOIN}}
\begin{columns}[c]
\begin{column}{0.3\textwidth}\centering
\joinvenn{Inner Join}{white}{joinfill}{white}
\end{column}
\begin{column}{0.65\textwidth}
\begin{lstlisting}[language=SQL]
SELECT
    students.*,
    teachers.Teacher_Name
FROM students INNER JOIN teachers
    ON students.Class = teachers.Class;
\end{lstlisting}
\end{column}
\end{columns}
\vspace{0.5em}
\centering
{\scriptsize\setlength{\tabcolsep}{4pt}
\begin{tabular}{@{}lcccl@{}}
\multicolumn{5}{@{}l}{\textbf{Inner Join}}\\ \toprule
Student\_Name & Grad\_Year & Class & Grade & Teacher\_Name \\ \midrule
Sandy & 11 & Gym & A & B. Joe \\
Gary & 9 & Civics & B+ & S. Small \\
Pearl & 12 & Calculus & B & N. Watterson \\
Karen & 12 & Calculus & C- & N. Watterson \\
Larry & 10 & Spelling & B & T. Wilson \\
Plankton & 11 & Gym & B- & B. Joe \\
Squidward & 10 & Biology & A- & C. Krueger \\
Bob & 11 & Biology & B+ & C. Krueger \\ \bottomrule
\end{tabular}}
\end{frame}

\begin{frame}[fragile]{Joining Tables: \code{LEFT JOIN}}
\begin{columns}[c]
\begin{column}{0.3\textwidth}\centering
\joinvenn{Left Join}{joinfill}{joinfill}{white}
\end{column}
\begin{column}{0.65\textwidth}
\begin{lstlisting}[language=SQL]
SELECT
    students.*,
    teachers.Teacher_Name
FROM students LEFT JOIN teachers
    ON students.Class = teachers.Class;
\end{lstlisting}
\end{column}
\end{columns}
\vspace{0.4em}
\centering
{\scriptsize\setlength{\tabcolsep}{4pt}
\begin{tabular}{@{}lcccl@{}}
\multicolumn{5}{@{}l}{\textbf{Left Join}}\\ \toprule
Student\_Name & Grad\_Year & Class & Grade & Teacher\_Name \\ \midrule
Patrick & 9 & Lunch & A+ & \red{(blank)} \\
Sandy & 11 & Gym & A & B. Joe \\
Gary & 9 & Civics & B+ & S. Small \\
Pearl & 12 & Calculus & B & N. Watterson \\
Karen & 12 & Calculus & C- & N. Watterson \\
Larry & 10 & Spelling & B & T. Wilson \\
Plankton & 11 & Gym & B- & B. Joe \\
Squidward & 10 & Biology & A- & C. Krueger \\
Bob & 11 & Biology & B+ & C. Krueger \\ \bottomrule
\end{tabular}}

\vspace{0.3em}
{\scriptsize\red{Note: there is no match in the Teachers table for Lunch.}}
\end{frame}

\begin{frame}[fragile]{Joining Tables: \code{FULL JOIN}}
\begin{columns}[c]
\begin{column}{0.3\textwidth}\centering
\joinvenn{Full Join}{joinfill}{joinfill}{joinfill}
\end{column}
\begin{column}{0.65\textwidth}
\begin{lstlisting}[language=SQL]
SELECT
    students.*,
    teachers.*
FROM students FULL JOIN teachers
    ON students.Class = teachers.Class;
\end{lstlisting}
\end{column}
\end{columns}
\vspace{0.2em}
\centering
{\tiny\setlength{\tabcolsep}{4pt}
\begin{tabular}{@{}lccccc@{}}
\multicolumn{6}{@{}l}{\textbf{Full Join}}\\ \toprule
Student\_Name & Grad\_Year & Class & Grade & Teacher\_Name & \red{Class1} \\ \midrule
Patrick & 9 & Lunch & A+ & & \\
Sandy & 11 & Gym & A & B. Joe & Gym \\
Gary & 9 & Civics & B+ & S. Small & Civics \\
Pearl & 12 & Calculus & B & N. Watterson & Calculus \\
Karen & 12 & Calculus & C- & N. Watterson & Calculus \\
Larry & 10 & Spelling & B & T. Wilson & Spelling \\
Plankton & 11 & Gym & B- & B. Joe & Gym \\
Squidward & 10 & Biology & A- & C. Krueger & Biology \\
Bob & 11 & Biology & B+ & C. Krueger & Biology \\
 & & & & R. Robinson & Drama \\
 & & & & B. Parham & Band \\
 & & & & G. Watterson & Music Theory \\ \bottomrule
\end{tabular}}

\vspace{0.2em}
{\scriptsize\red{Note: both tables have a \code{Class} column, so a new name is given in the
joined table: \code{Class1}. Also, unmatched items are given blank values in the joined table.}}
\end{frame}

\begin{frame}{Table Management: Practice}
\begin{prob}
\begin{enumerate}
  \item Using a \code{LEFT JOIN}, append \code{UrbanPop} and \code{Theft} from the
        \code{SCHEMA\_C.USArrests} table to the \code{SCHEMA\_A.Iris} table using \code{ID}.
  \vspace{0.5em}
  \item Using an \code{INNER JOIN}, combine the \code{SCHEMA\_C.USArrests} and
        \code{SCHEMA\_B.ToothGrowth} tables using \code{ID}.
\end{enumerate}
\end{prob}
\end{frame}

\begin{frame}[fragile]{Table Management: Practice -- \red{Answers}}
\small
1. Using a \code{LEFT JOIN}, append \code{UrbanPop} and \code{Theft} from
\code{SCHEMA\_C.USArrests} to \code{SCHEMA\_A.Iris} using \code{ID}.
\begin{lstlisting}[language=SQL]
SELECT IRIS.*, USARRESTS.URBANPOP, USARRESTS.THEFT
FROM SCHEMA_A.IRIS
LEFT JOIN SCHEMA_C.USARRESTS
    ON SCHEMA_A.IRIS.ID = SCHEMA_C.USARRESTS.ID
\end{lstlisting}
\vspace{0.4em}
2. Using an \code{INNER JOIN}, combine \code{SCHEMA\_C.USArrests} and
\code{SCHEMA\_B.ToothGrowth} using \code{ID}.
\begin{lstlisting}[language=SQL]
SELECT * FROM SCHEMA_C.USARRESTS
INNER JOIN SCHEMA_B.TOOTHGROWTH
    ON SCHEMA_C.USARRESTS.ID = SCHEMA_B.TOOTHGROWTH.ID
\end{lstlisting}
\end{frame}

\begin{frame}[fragile]{Note: Three Ways to Write an \code{INNER JOIN}}
In SQL, you can do an \textit{INNER JOIN} (answer \#2) in several ways:
\begin{lstlisting}[language=SQL]
SELECT * FROM SCHEMA_C.USARRESTS
    INNER JOIN SCHEMA_B.TOOTHGROWTH
    ON SCHEMA_C.USARRESTS.ID = SCHEMA_B.TOOTHGROWTH.ID
\end{lstlisting}
\textit{or}
\begin{lstlisting}[language=SQL]
SELECT * FROM SCHEMA_C.USARRESTS
    JOIN SCHEMA_B.TOOTHGROWTH
    ON SCHEMA_C.USARRESTS.ID = SCHEMA_B.TOOTHGROWTH.ID
\end{lstlisting}
\textit{or}
\begin{lstlisting}[language=SQL]
SELECT * FROM SCHEMA_C.USARRESTS,
              SCHEMA_B.TOOTHGROWTH
    WHERE SCHEMA_C.USARRESTS.ID = SCHEMA_B.TOOTHGROWTH.ID
\end{lstlisting}
\end{frame}

%======================================================================
\section{Aliasing and CASE WHEN Statements}
%======================================================================

\sectionframe

\begin{frame}[fragile]{Aliasing}
\begin{itemize}
  \item Reference tables or subqueries in code with a shortcut.
  \item Aliases allow you to refer to tables multiple times in your code with a
        shortcut, which gives you much cleaner code.
\end{itemize}
\begin{columns}[T]
\begin{column}{0.48\textwidth}
\textit{e.g.\ Instead of using}
\begin{lstlisting}[language=SQL,basicstyle=\ttfamily\tiny]
SELECT
    students.*,
    teachers.Teacher_Name
FROM students LEFT JOIN teachers
    ON students.Class = teachers.Class;
\end{lstlisting}
\end{column}
\begin{column}{0.48\textwidth}
\textit{We can use}
\begin{lstlisting}[language=SQL,basicstyle=\ttfamily\tiny]
SELECT
    s.*,
    t.Teacher_Name
FROM students s LEFT JOIN teachers t
    ON s.Class = t.Class;
\end{lstlisting}
\end{column}
\end{columns}
\end{frame}

\begin{frame}[fragile]{Aliasing De-clutters Your Code}
Aliasing can be \textbf{VERY} useful in de-cluttering your code!\\
{\small\red{NOTE: Use abbreviations that tie back to your table names! Try not to use
a, b, c, T1, T2, \dots\ to alias.}}

\textit{e.g.\ Instead of using}
\begin{lstlisting}[language=SQL,basicstyle=\ttfamily\tiny]
SELECT FAKE_SCHEMA.EXAMPLE_DATA_VW.POL_NAME,
FAKE_SCHEMA.EXAMPLE_DATA_VW.POL_EXPIRATION, FAKE_SCHEMA_TWO.EXAMPLE_DATA_TABLE.POL_NUMBER
FROM FAKE_SCHEMA.EXAMPLE_DATA_VW LEFT JOIN FAKE_SCHEMA_TWO.EXAMPLE_DATA_TABLE ON
FAKE_SCHEMA.EXAMPLE_DATA_VW.POL_KEY = FAKE_SCHEMA_TWO.EXAMPLE_DATA_TABLE.POL_KEY AND
FAKE_SCHEMA.EXAMPLE_DATA_VW.NEW_RENEWAL_CODE = FAKE_SCHEMA_TWO.EXAMPLE_DATA_TABLE.NEW_RENEWAL_CODE
\end{lstlisting}
\textit{We can use}
\begin{lstlisting}[language=SQL,basicstyle=\ttfamily\tiny]
SELECT EDV.POL_NAME, EDV.POL_EXPIRATION, EDT.POL_NUMBER
FROM FAKE_SCHEMA.EXAMPLE_DATA_VW EDV
     LEFT JOIN FAKE_SCHEMA_TWO.EXAMPLE_DATA_TABLE EDT
          ON EDV.POL_KEY = EDT.POL_KEY
          AND EDV.NEW_RENEWAL_CODE = EDT.NEW_RENEWAL_CODE
\end{lstlisting}
\end{frame}

\begin{frame}{Aliasing: Practice -- \textit{use aliasing!}}
\begin{prob}\small
\begin{enumerate}
  \item Using a \code{LEFT JOIN}, append \code{weight} from the \code{SCHEMA\_B.PlantGrowth}
        table to the \code{SCHEMA\_A.mtcars} table by matching \code{color\_ID} in \code{mtcars}
        to \code{ID} in \code{PlantGrowth}, selecting \code{Name}, \code{mpg}, \code{cyl}, and \code{wt} from the \code{mtcars} table.
        Also, create a new column called \code{weight\_diff} that gives the difference of the
        \code{weight} column and the \code{wt} column.
  \vspace{0.5em}
  \item Using \code{INNER JOIN}, combine the \code{SCHEMA\_A.Iris}, \code{SCHEMA\_B.PlantGrowth},
        and \code{SCHEMA\_C.USArrests} tables using \code{ID} for \code{iris} and
        \code{PlantGrowth}, and \code{State} for \code{iris} and \code{USArrests}.
        Select all of the data in the \code{iris} table, the \code{weight} column from the
        \code{PlantGrowth} table, and the \code{Murder}, \code{State}, and \code{Assault}
        columns from the \code{USArrests} table.
\end{enumerate}
\end{prob}
\end{frame}

\begin{frame}[fragile]{Aliasing: Practice -- \red{Answers} (1)}
{\small 1. \code{LEFT JOIN} \code{weight} from \code{SCHEMA\_B.PlantGrowth} onto
\code{SCHEMA\_A.mtcars} (\code{color\_ID} $\leftrightarrow$ \code{ID}), selecting \code{Name},
\code{mpg}, \code{cyl}, \code{wt}, and the new column \code{weight\_diff} = \code{weight} $-$ \code{wt}.}
\begin{lstlisting}[language=SQL]
SELECT
    M.NAME,
    M.MPG,
    M.CYL,
    M.WT,
    P.WEIGHT,
    (P.WEIGHT - M.WT) AS WEIGHT_DIFF
FROM SCHEMA_A.MTCARS M
        LEFT JOIN SCHEMA_B.PLANTGROWTH P
            ON M.COLOR_ID = P.ID
\end{lstlisting}
\end{frame}

\begin{frame}[fragile]{Aliasing: Practice -- \red{Answers} (2)}
{\small 2. \code{INNER JOIN} of \code{SCHEMA\_A.Iris}, \code{SCHEMA\_B.PlantGrowth} (on \code{ID})
and \code{SCHEMA\_C.USArrests} (on \code{State}): all of \code{iris}, \code{weight} from
\code{PlantGrowth}, and \code{Murder}, \code{State}, \code{Assault} from \code{USArrests}.}
\begin{lstlisting}[language=SQL]
SELECT
    I.*,
    P.WEIGHT,
    U.MURDER,
    U.STATE,
    U.ASSAULT
FROM SCHEMA_A.IRIS I
        INNER JOIN SCHEMA_B.PLANTGROWTH P
            ON I.ID = P.ID
        INNER JOIN SCHEMA_C.USARRESTS U
            ON I.STATE = U.STATE
\end{lstlisting}
\end{frame}

\begin{frame}[fragile]{String Handling (1)}
\begin{itemize}\small
  \item \lstinline[language=SQL]|SUBSTR(x, start [, length]);|
  \begin{itemize}
    \item Returns a substring of \code{x} that begins at the position specified by \code{start}.
          An optional \code{length} for the substring may be supplied.
    \item \code{start} can be negative to count from the end of the string.
  \end{itemize}
  {\scriptsize
  \begin{tabular}{@{}ll@{}}
    \textcolor{gray}{\textit{policy=88UECAB4758}} & \code{SUBSTR(policy,6,6)} = \textcolor{gray}{\textit{AB4758}} \\
    \textcolor{gray}{\textit{policy=88UECAB4758}} & \code{SUBSTR(policy,-4)} = \textcolor{gray}{\textit{4758}}
  \end{tabular}}
  \item \lstinline[language=SQL]|CONCAT(x, y);|
  \begin{itemize}
    \item Concatenates the strings \code{x} and \code{y} and returns the appended string.
    \item You can also use \code{||} to concatenate strings.
  \end{itemize}
  {\scriptsize
  \begin{tabular}{@{}ll@{}}
    \textcolor{gray}{\textit{pol\_sym=UEC}} & \code{CONCAT(pol\_sym,pol\_num)} = \textcolor{gray}{\textit{UECAB4758}} \quad \textit{or} \\
    \textcolor{gray}{\textit{pol\_num=AB4758}} & \code{(pol\_sym || pol\_num)} = \textcolor{gray}{\textit{UECAB4758}}
  \end{tabular}}
  \item \lstinline[language=SQL]|LENGTH(x);|
  \begin{itemize}
    \item Returns the number of characters in \code{x}.
  \end{itemize}
  {\scriptsize
  \begin{tabular}{@{}ll@{}}
    \textcolor{gray}{\textit{policy=88UECAB4758}} & \code{LENGTH(policy)} = \textcolor{gray}{\textit{11}}
  \end{tabular}}
\end{itemize}
\end{frame}

\begin{frame}[fragile]{String Handling (2)}
\begin{itemize}
  \item \lstinline[language=SQL]|NVL(x, value);|
  \begin{itemize}
    \item Returns \code{value} if \code{x} is null; otherwise, \code{x} is returned.
          \quad e.g.\ \code{NVL(prem,0)}
  \end{itemize}
  \vspace{0.4em}
  \item \lstinline[language=SQL]|LTRIM(x [, trim_string]);| / \lstinline[language=SQL]|RTRIM(x [, trim_string]);|
  \begin{itemize}
    \item Trims characters from the left or right of \code{x}
          (by default, leading/trailing spaces).
  \end{itemize}
  \vspace{0.4em}
  \item \lstinline[language=SQL]|CHARINDEX(search_string, x);|
  \begin{itemize}
    \item Searches \code{x} for \code{search\_string} and returns the position of its first
          occurrence (0 if not found). Note the order: the string to find comes \emph{first}.
    \item In Oracle/SQLite the equivalent is \lstinline[language=SQL]|INSTR(x, search_string)|.
  \end{itemize}
  \vspace{0.4em}
  \item \lstinline[language=SQL]|REPLACE(x, search_string, replace_string);|
  \begin{itemize}
    \item Searches \code{x} for \code{search\_string} and replaces it with \code{replace\_string}.
  \end{itemize}
\end{itemize}
\end{frame}

\begin{frame}{String Handling: Practice}
\begin{prob}
Write a query that extracts two columns from the \code{SCHEMA\_A.mtcars} table, which
decomposes the \code{Name} into two parts (separated by the first space): \code{Brand}
and \code{Model}.
\end{prob}
\vspace{0.5em}
\begin{prob}
Write a query that extracts the accident year, month, and day from \code{acc\_dt} of the
\code{SCHEMA\_D.lossexpense} table. Further, create a column with date format using the
extracted information.
\end{prob}
\vspace{0.5em}
\begin{prob}
Using a \code{LEFT JOIN}, append all columns from the \code{SCHEMA\_D.premexposures}
table to the \code{SCHEMA\_D.lossexpense} table using \code{pol\_id}. Note that
\code{pol\_id} can be obtained in the \code{SCHEMA\_D.lossexpense} table by concatenating
\code{pol\_regsym} and \code{pol\_num}.
\end{prob}
\end{frame}

\begin{frame}[fragile]{String Handling: Practice -- \red{Answers}}
{\small Split \code{Name} of \code{mtcars} into \code{Brand} and \code{Model}:}
\begin{lstlisting}[language=SQL]
SELECT mtcars.Name,
    SUBSTR(mtcars.Name, 1, CHARINDEX(' ', mtcars.Name)-1) AS Brand,
    SUBSTR(mtcars.Name, CHARINDEX(' ', mtcars.Name)+1,
           LENGTH(mtcars.Name)) AS Model
FROM mtcars;
\end{lstlisting}
\vspace{0.3em}
{\small Append \code{premexposures} to \code{lossexpense} via the concatenated \code{pol\_id}:}
\begin{lstlisting}[language=SQL]
SELECT le.*, pe.*
FROM SCHEMA_D.lossexpense le LEFT JOIN SCHEMA_D.premexposures pe
    ON CONCAT(le.pol_regsym, le.pol_num) = pe.pol_id;
\end{lstlisting}
\end{frame}

\begin{frame}[fragile]{String Handling: Practice -- \red{Answers} (cont.)}
{\small Extract accident year, month, and day from \code{acc\_dt} (format \code{m/d/yyyy}),
then build a date column from the extracted parts:}
\begin{lstlisting}[language=SQL,basicstyle=\ttfamily\fontsize{5.6}{6.6}\selectfont]
SELECT le.acc_dt,
    SUBSTR(le.acc_dt, -4) AS acc_year,
    SUBSTR(le.acc_dt, 1, CHARINDEX('/', le.acc_dt)-1) AS acc_month,
    SUBSTR(SUBSTR(le.acc_dt, CHARINDEX('/', le.acc_dt)+1), 1,
           CHARINDEX('/', SUBSTR(le.acc_dt, CHARINDEX('/', le.acc_dt)+1))-1) AS acc_day,
    TO_DATE(SUBSTR(le.acc_dt, -4) || '-' ||
            SUBSTR(le.acc_dt, 1, CHARINDEX('/', le.acc_dt)-1) || '-' ||
            SUBSTR(SUBSTR(le.acc_dt, CHARINDEX('/', le.acc_dt)+1), 1,
                   CHARINDEX('/', SUBSTR(le.acc_dt, CHARINDEX('/', le.acc_dt)+1))-1),
            'YYYY-MM-DD') AS acc_date
FROM SCHEMA_D.lossexpense le;
\end{lstlisting}
\begin{itemize}\small
  \item \code{TO\_DATE(string, format)} converts text to a date; \code{'YYYY-MM-DD'} tells it how
        the pieces are laid out.
  \item Shortcut once you know the format: \code{TO\_DATE(le.acc\_dt, 'MM/DD/YYYY')}.
\end{itemize}
\end{frame}

\begin{frame}[fragile]{\code{CASE} Statement}
\begin{itemize}
  \item Effectively an IF-ELSE statement.
  \item \code{CASE} looks for the first match; once it meets this match, it performs the
        \code{THEN} and goes to the next row of the data.
\end{itemize}
\begin{lstlisting}[language=SQL]
SELECT c.*,
       CASE
          WHEN c.city = 'Seattle' THEN 'Northwest'
          WHEN c.city = 'Orlando' THEN 'Southeast'
          ELSE 'Merica'
       END AS REGION
  FROM customers c
\end{lstlisting}
{\small\color{red!70!black}
\begin{itemize}
  \item Use an \code{ELSE} statement when you want the remaining circumstances to be
        something other than blank.
  \item Don't forget the \code{END} statement!
  \item \code{AS REGION} names the column being created.
\end{itemize}}
\end{frame}

\begin{frame}[fragile]{\code{CASE} Statement: Practice}
\begin{prob}\small
Write a query that creates an additional column in the \code{SCHEMA\_C.USArrests} table
using a \code{CASE} statement to give unique codes based on the values in the \code{Murder}
column. For example, if the \code{Murder} rate is greater than 10, the new column could be
`15+' or `High' or `Terrifying!'.
\end{prob}
\pause
\textbf{\red{Answer:}}
\begin{lstlisting}[language=SQL]
SELECT
    U.*,
    CASE WHEN U.MURDER > 15 THEN 'AHHHHH!!!!'
         WHEN U.MURDER > 10 THEN 'SCARY AT NIGHT!'
         WHEN U.MURDER > 5 THEN 'JUST BE CAREFUL'
         ELSE 'PEACE, LOVE, AND HAPPINESS'
    END AS MURDER_CODE
FROM SCHEMA_C.USARRESTS U
\end{lstlisting}
\end{frame}

%======================================================================
\section{Nested Queries and WITH Statements}
%======================================================================

\sectionframe

\begin{frame}[fragile]{Subqueries: Motivation}
\begin{itemize}\small
  \item Oftentimes, you will need more than one table in more than one step to get what you need.
  \item What if you want to accomplish this without creating multiple tables in the
        database? Or\dots what if you want to accomplish this without creating mid-step tables?
\end{itemize}
For example:
\begin{lstlisting}[language=SQL]
CREATE TABLE Average_Age AS
SELECT
    AVG(AGE) AS AvAge
FROM SCHEMA.TABLE;

CREATE TABLE Greater_than_Average AS
SELECT
    NAMES
FROM SCHEMA.TABLE nm, Average_Age aa
WHERE nm.AGE > aa.AvAge;
\end{lstlisting}
\end{frame}

\begin{frame}[fragile]{Subquery}
\begin{itemize}
  \item SQL query nested in a larger query
  \begin{itemize}
    \item Can be added to the \textit{WHERE}, \textit{SELECT}, and \textit{FROM} clauses
    \item Can be run on their own to produce results, often single or multiple values
  \end{itemize}
  \item Queries with subqueries execute from the innermost query outward
  \item Multiple nested subqueries are possible
\end{itemize}
\vspace{0.5em}
\centering
\begin{tikzpicture}
  \node[draw=dorange, very thick, rounded corners, inner sep=4pt] (outer) {%
    \begin{minipage}{0.8\textwidth}
\begin{lstlisting}[language=SQL,frame=none,backgroundcolor={}]
SELECT NAMES
  FROM SCHEMA.TABLE
  WHERE AGE > (SELECT AVG(AGE) FROM SCHEMA.TABLE);
\end{lstlisting}
    \end{minipage}};
  \node[font=\small\bfseries, text=dorange, below left=0.35cm and 0.6cm of outer.south] (lo) {Outer Query};
  \node[font=\small\bfseries, text=dgreen, below right=0.35cm and -0.6cm of outer.south] (li) {Inner Query};
  \draw[->, dorange, thick] (lo.north) -- ++(0.5,0.35);
  \draw[->, dgreen, thick] (li.north) -- ++(0.3,0.75);
\end{tikzpicture}
\end{frame}

\begin{frame}[fragile]{Subquery: Practice}
\begin{prob}
Using a subquery, return all rows from the \code{iris} table with \code{ID}s the same as
the \code{color\_ID}s associated with Mercedes found in the \code{mtcars} table.
\textit{HINT: the wildcard character in SQL is \%.}
\end{prob}
\pause
\vspace{0.5em}
\textbf{\red{Answer:}}
\begin{lstlisting}[language=SQL]
SELECT * FROM IRIS
    WHERE ID IN (SELECT COLOR_ID FROM MTCARS
                     WHERE NAME LIKE 'Merc%')
\end{lstlisting}
\end{frame}

\begin{frame}[fragile]{Subquery: Scoping Considerations}
\begin{itemize}\small
  \item Be mindful of the order of operations; the \underline{scope} of any query is
        contained in the \textbf{\textit{FROM}} clause.
  \item Any subquery in the \textit{WHERE} clause of an outer query will return values
        and not a table that can be selected from.
\end{itemize}
\textit{For example:} \red{\small\textit{we cannot pull \code{AVG\_AGE} in the outer query
from this inner query}} \textcolor{red}{$\times$}
\begin{lstlisting}[language=SQL]
SELECT NAMES, AVG_AGE
FROM SCHEMA.TABLE
WHERE AGE > (SELECT AVG(AGE) AS AVG_AGE FROM SCHEMA.TABLE);
\end{lstlisting}
\red{\small\textit{This code will work correctly}} \textcolor{dgreen}{$\checkmark$}
\begin{lstlisting}[language=SQL]
SELECT sch.NAMES, iq.AVG_AGE
FROM SCHEMA.TABLE sch,
     (SELECT AVG(AGE) AS AVG_AGE FROM SCHEMA.TABLE) iq
WHERE sch.AGE > iq.AVG_AGE;
\end{lstlisting}
\end{frame}

\begin{frame}[fragile]{Subquery: \code{WITH}}
\begin{itemize}\small
  \item Special subquery that generates a table to be used within a larger query.
  \item Very useful for organizing complex queries (\textbf{MUCH} easier to read than nested
        queries, particularly when there are multiple subqueries!)
\end{itemize}
\textit{For example:}
\begin{lstlisting}[language=SQL]
WITH temp AS (SELECT AVG(AGE) AS AVG_AGE FROM SCHEMA.TABLE)
SELECT sch.NAMES, iq.AVG_AGE
FROM SCHEMA.TABLE sch, temp iq
WHERE sch.AGE > iq.AVG_AGE;
\end{lstlisting}
\begin{lstlisting}[language=SQL]
SELECT sch.NAMES, iq.AVG_AGE
FROM SCHEMA.TABLE sch,
     (SELECT AVG(AGE) AS AVG_AGE FROM SCHEMA.TABLE) iq
WHERE sch.AGE > iq.AVG_AGE;
\end{lstlisting}
\red{\small\textit{Both result in the same output, but the top one is easier to read.}}
\end{frame}

\begin{frame}[fragile]{\code{WITH}: Practice}
\begin{prob}
Using a \code{WITH} statement, return all rows from the \code{iris} table with \code{ID}s
the same as the \code{color\_ID}s associated with Mercedes found in the \code{mtcars} table.
\textit{HINT: the wildcard character in SQL is \%.}
\end{prob}
\pause
\vspace{0.5em}
\textbf{\red{Answer:}}
\begin{lstlisting}[language=SQL]
WITH TEMP AS
(SELECT COLOR_ID FROM MTCARS
        WHERE NAME LIKE 'Merc%')
SELECT I.* FROM IRIS I, TEMP T
    WHERE I.ID = T.COLOR_ID
\end{lstlisting}
\end{frame}

\begin{frame}{Subquery and \code{WITH}: Practice}
\begin{prob}
Return all rows from the \code{ToothGrowth} table along with the \code{Murder} column from
the \code{USArrests} table for urban populations greater than 80.

\vspace{0.5em}
Do this using a subquery, and again using a \code{WITH} statement.
\end{prob}
\end{frame}

\begin{frame}[fragile]{Subquery and \code{WITH}: Practice -- \red{Answers}}
\textbf{Using a subquery:}
\begin{lstlisting}[language=SQL]
SELECT
    T.*, P.MURDER
FROM TOOTHGROWTH T
INNER JOIN (SELECT U.*
            FROM USARRESTS U
            WHERE U.URBANPOP > 80) P
    ON T.ID = P.ID
\end{lstlisting}
\vspace{0.3em}
\textbf{Using a \code{WITH} statement:}
\begin{lstlisting}[language=SQL]
WITH POP AS (SELECT U.*
             FROM USARRESTS U WHERE U.URBANPOP > 80)
SELECT
    T.*, P.MURDER
FROM TOOTHGROWTH T
INNER JOIN POP P
    ON T.ID = P.ID
\end{lstlisting}
\end{frame}

\end{document}
